\section{A continuous weight with one zero}
\label{sec:continuity}

The measurable obstruction of Theorem~\ref{thm:measurable} is now in place.
We will replace its logarithmic cost $g_*$ by a function continuous off one
point and tending to infinity at that point.  The replacement changes the
cost by the constant $1$ plus a measurable coboundary.  Since its transfer function need not be
bounded, we begin with the precise operator statement that will preserve
the obstruction.

\subsection{Transport by an affiliated positive operator}

A positive self-adjoint operator is \emph{affiliated} with a von Neumann
algebra if all its spectral projections belong to that algebra.  Such an
operator commutes, including on its domain, with every unitary in the
commutant.
The action of an invertible affiliated operator on projection lattices
also appears in \citet[Proposition~3.1]{mori}.  The form needed here
includes a specified intertwining domain and preservation of the trace;
we give a direct spectral-cutoff proof.

\begin{lemma}\label{lem:affiliated-transport}
Let $(\mathcal M,\tau)$ be a finite von Neumann algebra acting on a Hilbert
space, with faithful normal tracial state $\tau$.  Let $H$ be an injective
positive self-adjoint operator affiliated with $\mathcal M$, and let
$A,B\in\mathcal M$.  Suppose
\[
 A(\operatorname{Dom}H)\subseteq\operatorname{Dom}H,
 \qquad HA\zeta=BH\zeta\quad(\zeta\in\operatorname{Dom}H).
\]
For every projection $p\in\mathcal M$ with $A\operatorname{ran}p\subseteq
\operatorname{ran}p$, the projection $P$ onto
\[
 \overline{H(\operatorname{ran}p\cap\operatorname{Dom}H)}
\]
belongs to $\mathcal M$, is invariant under $B$, and satisfies
$\tau(P)=\tau(p)$.
\end{lemma}

\begin{proof}
Every commutant unitary preserves $\operatorname{ran}p$ and
$\operatorname{Dom}H$, and commutes with $H$ on that domain.  It therefore
maps the displayed closed space onto itself.  Its projection commutes with
the commutant, proving $P\in\mathcal M$.  If
$\zeta\in\operatorname{ran}p\cap\operatorname{Dom}H$, both invariance
assumptions give
$A\zeta\in\operatorname{ran}p\cap\operatorname{Dom}H$, and
$BH\zeta=HA\zeta$.  Thus $B$ preserves the unclosed image and, by
boundedness, its closure.

It remains to prove the trace assertion; bounded similarity cannot be
assumed.  Put $E_N=\mathbf1_{[1/N,N]}(H)$.  Since $H$ is injective,
$E_N\uparrow1$.  For $q_N=p\wedge E_N$, the bounded element $Hq_N=(HE_N)q_N$ of
$\mathcal M$ is bounded below by $1/N$ on its initial space.  Its image is
closed, and its polar decomposition identifies the initial projection
$q_N$ with its final projection, which is dominated by $P$.  Therefore
\[
 \tau(P)\ge\tau(p\wedge E_N)
 \ge\tau(p)+\tau(E_N)-1.
\]
The last inequality follows directly from the polar decomposition of
$(1-E_N)p$: its initial support is $p-(p\wedge E_N)$ and its final support
is at most $1-E_N$.  Letting $N$ increase gives $\tau(P)\ge\tau(p)$.

For the reverse inequality, set $q'_N=(1-p)\wedge E_N$.  The bounded
operator $H^{-1}q'_N=(H^{-1}E_N)q'_N$ is likewise bounded below on its initial space and
has closed image.  This image is orthogonal to $\operatorname{ran}P$.
Indeed, for $\eta\in\operatorname{ran}q'_N$ and
$\zeta\in\operatorname{ran}p\cap\operatorname{Dom}H$, self-adjointness
gives
\[
 \langle H^{-1}\eta,H\zeta\rangle=\langle\eta,\zeta\rangle=0;
\]
here $H^{-1}\eta$ lies in $\operatorname{Dom}H$.  Polar decomposition and
the same meet estimate yield
\[
 \tau(1-P)\ge\tau((1-p)\wedge E_N)
 \ge\tau(1-p)+\tau(E_N)-1.
\]
Passing to the limit proves the opposite inequality for $\tau(P)$.
\end{proof}

\subsection{Smoothing and moving the logarithmic drops}

We claim that there exist a nonnegative function $g$, finite and continuous
on $X\setminus\{0\}$ and tending to infinity at $0$, and an almost
everywhere finite real measurable function $b$ such that
\begin{equation}\label{eq:continuous-coboundary}
 g=g_*+1+b\circ\sigma-b\quad\text{almost everywhere},
 \qquad \int_Xg\,d\mu=\infty.
\end{equation}
The construction has three parts: smooth each drop, move its support into
a shrinking arc, and add an integrable background cost that guarantees
divergence at the center of those arcs.  Apart from the constant $1$, all changes will be explicit
coboundaries.

Fix $i\in J_j$ and put $\delta_i=\alpha_i/4$.  For $u\in X$, define the
continuous periodic smoothing of the sawtooth by
\[
 s_i(u)=\frac1{2\delta_i}\int_{-\delta_i}^{\delta_i}\{u+t\}\,dt.
\]
It equals $\{u\}$ outside the $\delta_i$-neighborhood of the cut at $0$.
Set
\[
 h_i(x)=C_j\bigl(\alpha_i+s_i(u_i(x))
                         -s_i(u_i(x)+\alpha_i)\bigr).
\]
The elementary identity
$\lfloor\{u\}+\alpha_i\rfloor=
\alpha_i+\{u\}-\{u+\alpha_i\}$ shows that $h_i$ is the corresponding
average of the original one-coordinate cost.  It is nonnegative and
continuous, with
\begin{equation}\label{eq:smoothed-support}
 \int_X h_i\,d\mu=C_j\alpha_i,
 \qquad \mu(\operatorname{supp}h_i)\le\alpha_i+2\delta_i.
\end{equation}
For
\[
 d_i^{\rm sm}(x)=C_j\bigl(s_i(u_i(x))-u_i(x)\bigr),
\]
we have the exact measurable identity
\begin{equation}\label{eq:smoothing-coboundary}
 h_i(x)=C_j\lfloor u_i(x)+\alpha_i\rfloor
                  +d_i^{\rm sm}(x)-d_i^{\rm sm}(\sigma x).
\end{equation}
The support of $d_i^{\rm sm}$ has measure at most $2\delta_i$.

Recall the shrinking arcs $O_j$, the integers $h_j^*$, and the continuous
partition of unity $\phi_{j,n}$, $0\le n\le h_j^*$, prepared in the
parameter construction, with
$\sigma^n(\operatorname{supp}\phi_{j,n})\subset O_j$.  Define
\[
 h'_i(x)=\sum_{n=0}^{h_j^*}(h_i\phi_{j,n})(\sigma^{-n}x).
\]
This is continuous and nonnegative, supported in $O_j$, and has the same
integral as $h_i$.  Telescoping gives
\begin{equation}\label{eq:relocation-coboundary}
 h'_i=h_i+d'_i-d'_i\circ\sigma,
 \qquad
 d'_i(x)=\sum_{n=1}^{h_j^*}\sum_{r=1}^n
                      (h_i\phi_{j,n})(\sigma^{-r}x).
\end{equation}
Although this transfer contains a double sum, its support is contained in
just $h_j^*$ translates of $\operatorname{supp}h_i$, one for each possible
$r$.  Consequently
\[
 \sum_i\left(\mu\{d_i^{\rm sm}\ne0\}+\mu\{d'_i\ne0\}\right)
 \le C\sum_j\frac{1+h_j^*}{K_j}<\infty.
\]
Thus each of the transfer sums has only finitely many nonzero terms
almost everywhere, including at the translated point $\sigma x$.

The sum $\sum_i h'_i$ is locally finite on $X\setminus\{0\}$ because its
successive blocks are supported in $O_j\downarrow\{0\}$.  It is therefore
continuous there.  Its integral is infinite: Tonelli's theorem and
$\alpha_i>1/(m_jK_j)$ give
\begin{equation}\label{eq:infinite-cost-integral}
 \int_X\sum_i h'_i\,d\mu
 =\sum_j C_j\sum_{i\in J_j}\alpha_i
 \ge\sum_j C_j/K_j=\infty,
\end{equation}
since $C_j\ge K_j^2$.  Local finiteness alone need not force divergence
at $0$; the following background supplies that remaining property.

\subsection{A background cost diverging at the zero}

Let $W_j=2^{j+1}$.  Choose continuous functions $0\le v_j\le1$ with
$v_j(0)=1$, supported in closed arcs $I_j$ shrinking to $\{0\}$.  Choose
these arcs small enough that $\sigma I_j,\ldots,\sigma^{W_j}I_j$ are
pairwise disjoint and
\[
 \sum_j W_j\mu(I_j)<\infty.
\]
This is possible because each finite set of distinct points
$\theta,2\theta,\ldots,W_j\theta$ has positive minimum separation.
For $x\ne0$, put
\[
 g_0(x)=1+\sum_j\left(v_j(x)-\frac1{W_j}
                              \sum_{n=1}^{W_j}v_j(\sigma^{-n}x)\right).
\]
For each $j$, disjointness bounds the norm of the subtracted function by
$1/W_j$.  These bounds sum to $1/2$, so the negative series converges
uniformly.  The positive series is locally finite away from $0$.
Consequently $g_0$ is continuous there and $g_0\ge1/2$.  For every fixed
integer $J$, the finite sum $\sum_{j\le J}v_j(x)$ tends to $J$ as $x\to0$.
The bounded negative series therefore implies $g_0(x)\to\infty$ as
$x\to0$.  Its integral is finite, since
$\sum_j\int v_j\,d\mu\le\sum_j\mu(I_j)<\infty$ and the negative series
is uniformly bounded.

Define
\[
 b_0(x)=\sum_j\frac1{W_j}\sum_{n=1}^{W_j}\sum_{r=1}^n
                                      v_j(\sigma^{-r}x).
\]
The $j$th summand is supported in $\bigcup_{r=1}^{W_j}\sigma^r I_j$.
The assumed summability makes this sum finite almost everywhere.  On a
conull set where the sums at $x$ and $\sigma x$ are finite, telescoping
shows
\[
 g_0=1+b_0\circ\sigma-b_0.
\]
Now let
\[
 g(x)=g_0(x)+\sum_i h'_i(x)\quad(x\ne0),\qquad g(0)=\infty,
 \qquad b=b_0-\sum_i(d_i^{\rm sm}+d'_i).
\]
The preceding arguments prove that $g$ is nonnegative, finite and
continuous off $0$, tends to infinity at $0$, and has infinite integral.
The transfer $b$ is real and finite almost everywhere.  Summing
\eqref{eq:smoothing-coboundary} and \eqref{eq:relocation-coboundary} is
legitimate on the common conull set where the transfer sums and the
original cost have finitely many nonzero terms.  Together with the
background identity this proves \eqref{eq:continuous-coboundary}.

\subsection{Conclusion and spectral consequences}

Define the function on the unit circle by
\[
 f(e^{2\pi i x})=e^{-g(x)},\qquad e^{-\infty}=0.
\]
It is continuous, is positive away from $1$, and vanishes at $1$.
Its modulus is nonconstant and its zero set has Haar measure zero.  Moreover,
\begin{equation}\label{eq:zero-log-integral}
 \int_{\mathbb T}\log f\,dm=-\infty.
\end{equation}
In particular, its Fuglede--Kadison determinant is zero
\cite{fuglede-kadison}: the spectral distribution of $\pi(f)$ is the
Haar distribution of $f$, so the determinant definition is precisely the
exponential of \eqref{eq:zero-log-integral}.  Also $|U\pi(f)|=\pi(f)$,
which gives the same determinant for $T_f=U\pi(f)$.  This conclusion uses
the nonnegative integral calculation \eqref{eq:infinite-cost-integral};
it does not require the transfer $b$ to be integrable.

Choose a finite Borel version of $b$ and discard all integer translates
of the null set where the coboundary identity fails.
Let $H_b$ be multiplication on the $k$th fiber coordinate by
$e^{b(\sigma^k x)}$, with its natural square-integrability domain.  It is
positive, self-adjoint, and injective.  Its spectral projections are
$\pi$-images of spectral indicators of $e^b$, hence lie in $M$.  Thus
$H_b$ is affiliated with $M$.  The coboundary identity gives the pointwise
factor equality
\[
 e^{b(\sigma^{k+1}x)}e^{-g(\sigma^k x)}
       =e^{-1}f_*(\sigma^k x)e^{b(\sigma^k x)}.
\]
For every $\zeta\in\operatorname{Dom}H_b$, the right-hand side defines
the $k+1$ coordinate of $e^{-1}T_*H_b\zeta$, a square-integrable vector
because $T_*$ is bounded.  Therefore
\[
 T_f\zeta\in\operatorname{Dom}H_b,
 \qquad H_bT_f\zeta=e^{-1}T_*H_b\zeta.
\]
Lemma~\ref{lem:affiliated-transport} sends any invariant projection of
$T_f$ to an invariant projection of $e^{-1}T_*$ having the same trace.
Multiplication by the nonzero scalar $e^{-1}$ does not change invariant
ranges.  Theorem~\ref{thm:measurable} therefore excludes traces strictly
between $0$ and $1$.  This proves Theorem~\ref{thm:main} for the arbitrary
irrational $\theta$ fixed at the start.

The same continuous weight has a stronger spectral property, which we
record with an elementary proof.  This is the zero-determinant case of
\citet[Theorem~2.3]{zfs}.  We include the truncated-logarithm argument,
which uses uniform averaging under irrational rotation.

\begin{corollary}\label{cor:quasinilpotent}
The constructed operator $T_f$ is nonzero and quasinilpotent.
\end{corollary}

\begin{proof}
For $M>0$, the function
$\ell_M(x)=\log\max(f(e^{2\pi i x}),e^{-M})$ is continuous.
Its integrals decrease to $-\infty$ by
\eqref{eq:zero-log-integral}.  Irrational rotation averages of a continuous
function converge uniformly to its integral: this follows for Fourier
monomials from their geometric sums and then for continuous functions by
uniform trigonometric approximation.  Hence, for any $a>0$, first choosing
$M$ with $\int_X\ell_M\,d\mu<-2a$ and then taking $n$ sufficiently large
gives
\[
 \sup_{x\in X}\prod_{k=0}^{n-1}f(e^{2\pi i\sigma^k x})
 \le\sup_{x\in X}\exp\left(\sum_{k=0}^{n-1}\ell_M(\sigma^k x)\right)
 \le e^{-an}.
\]
The expression on the left is $\|T_f^n\|$, because the power is a unitary
$U^n$ times multiplication by that continuous product.  The spectral-radius
formula yields radius zero, since $a$ was arbitrary.  Finally $T_f\ne0$
because $f>0$ almost everywhere.
\end{proof}
